\chapter{Introducción}

El problema de la distribución eficiente de bienes en entornos urbanos es uno de los desafíos logísticos más relevantes tanto en la industria como en el campo del diseño algorítmico. En ciudades como el Cusco, cuya red vial histórica presenta calles angostas, sentidos únicos, zonas peatonales y calles frecuentemente bloqueadas por actividades culturales o de mantenimiento, la planificación manual de rutas de reparto resulta especialmente compleja e ineficiente.

El presente proyecto, desarrollado en el marco del curso de Programación III de la Escuela Profesional de Ingeniería Informática y de Sistemas (EPIIS) de la Universidad Nacional de San Antonio Abad del Cusco (UNSAAC), aborda esta problemática mediante la construcción de un sistema de gestión de rutas óptimas implementado como aplicación de escritorio en Python. El sistema integra una interfaz gráfica de usuario (GUI) basada en Tkinter que permite visualizar la red vial de Cusco ---20 lugares de interés unidos por calles reales de OpenStreetMap---, cerrar calles, gestionar pedidos de reparto y ejecutar cuatro paradigmas algorítmicos fundamentales del diseño de algoritmos: \textit{Greedy}, \textit{Divide y Vencerás}, \textit{Programación Dinámica} y \textit{Backtracking}.

El trabajo no se limita a la implementación funcional de los algoritmos, sino que incluye un análisis riguroso de su complejidad computacional temporal y espacial en notación Big-O, así como la ejecución de \textit{benchmarks} empíricos que permiten comparar el rendimiento real de cada paradigma ante diferentes volúmenes de datos. De este modo, el sistema sirve simultáneamente como herramienta práctica de apoyo a la logística urbana y como plataforma de demostración pedagógica de los paradigmas algorítmicos estudiados durante el semestre.

El documento se organiza de la siguiente manera: el Capítulo 2 describe el planteamiento del problema y su justificación; el Capítulo 3 establece los objetivos del proyecto; el Capítulo 4 desarrolla el marco teórico de los algoritmos implementados; el Capítulo 5 formula la hipótesis de trabajo; el Capítulo 6 describe la metodología de desarrollo; el Capítulo 7 presenta los resultados obtenidos y su discusión comparativa; y el Capítulo 8 expone las conclusiones y recomendaciones finales.
